\section{Related literature}

Technological change is a constant force of nature that has been around for thousands of years. The still ongoing debate referring to the effect on wages and employment has been parallel to almost every major techonlogical change. Recently, the introduction of \glsxtrfull{AI} and \glsxtrfull{LLM}'s reignited this debate since the capabilities shown by these technologies are extraordinary in the sense that they are capable of automating certain non routine cognitive tasks.The ambiguity of the effect on employment and wages that recent and historical studies have found inspires the scope of my main research question. How does the release of ChatGPT affect employment, wages and informality in Mexico? In order to understand this question I review literature surrounding technological change and its role in task automation, wages and employment. Second, I take a look at recent studies on the effect of \glsxtrshort{LLM}'s on employment and wages in an effort to understand the current state of the debate even though the literature is still new. Finally, the setting of my research question revolves around Mexico. Because of this, an important strand of literature I explore is the one that studies the dynamics of technological change in developing countries, escpecially Latin America.

Back in the 90's, there was a strand of literature trying to make sense of the increasing wage premium of skilled workers and technological unemployment. First, \cite{katz1992changes} develop a supply and demand model in an effort to understand the rising wage inequality and wage struture in the US. The authors find that the rising supply and demand for skilled workers shifted the college wage premium. They also attribute the changing wage structure to compositional shifts in relative labor demand within sectors due to skill biased technological change.  Then, \cite{krueger} examines employees who use a computer at work in an empirical framework. The author's estimates suggest a 10 to 15\% higher wage rate and above average wage growth in occupations with growing computer use. Meanwhile, \cite{cox}'s approach shows a temporary unemployment effect due to technological change. Throughout the decade some critics were issued \citep{dinardo} and by the end of the decade and the following years, \cite{katz1999} \& \cite{goldin2007} solidify the analysis by understanding the link between market forces and market institutions. The conclusion puts computers and fluctuations in the supply of more educated workers as the main suspects behind the skill wage differential.

What followed was a way to understand the mechanisms behind the skill premium after a technological change. \cite{murnane} investigate how the demand shift favored college labor. They come out with a framework capable of analyzing how computers change the demand of tasks on the job. They find that across industries, occupations and education groups, computers reduce routine manual and routine cognitive tasks and increase non routine cognitive tasks in labor input. Using this framework, research revolved around \textit{job polarization}. A phenomenon in which jobs are polarized into high wage and low wage, pushing out middle skill jobs  \citep{goos,autorkatzkearney}. The phenomenon is then linked with the growth of low skill jobs and falling costs of automating routine job tasks \citep{goos2014,autor2014}. Trying to understand the forces behind job polarization, \cite{acemoglu2018} explains two effects relating to the role of automation in task creation and destruction. The first one of these is a \textit{displacement effect} in which automation destroys tasks previously performed by labor. The second one is a \textit{productivity effect} that creates new tasks which are not yet able to be automated. These opposing effects result in ambiguous effects on wages and employment. The importance behind this framework relies on its capacity to explain how tasks exposed to automation are key in explaining the final outcome relating wages and employment. Tasks are not randomly exposed to automation. A technological change may involve destruction of tasks mostly performed by low or high skill workers and many authors exhibit concerns referring to wage inequality \citep{murnane,goldin2007,goos,goos2014,autor2014,brynjolfsson2014second,stiglitz2017artificial}. Because of the great importance behind this framework, I use it to disentangle and interpret the mechanisms behind my main results.

Branching out from this literature there is a new set of studies analyzing \glsxtrshort{AI}. The main difference this technology has showed is that \glsxtrshort{LLM}'s are capable of automating certain non routine cognitive tasks \citep{eloundou2023gptsgptsearlylook}. Although this difference shifts every notion of automated tasks we preiously had with computers, the framework established by \cite{acemoglu2018} is able to explain many details regarding the new technological paradigm. \cite{eloundou2023gptsgptsearlylook} develop an exposure index to \glsxtrfull{AI} by classifying 19,265 tasks from O*NET\footnote{A database that enlists occupation specific task statements in order to describe what a worker does in their respective occupation} with human tests and ChatGPT. The results point to 15\% of all tasks in the United States being completed much faster with the same quality when exposed to an \glsxtrshort{LLM}. This estimate increases to between 47\% and 65\% if additonal tools are considered to work alongisde the \glsxtrshort{LLM}. These findings point to a low share of tasks that can be fully automated. Moreover, \cite{eloundou2023gptsgptsearlylook} highlight higher income jobs facing higher exposure. The finding is further solidified by \cite{Felten2023OccupationalHI}, who characterizes exposed occupations as highly educated, highly paid and white collar.

Several empirical and behavioural papers have followed \cite{eloundou2023gptsgptsearlylook} trying to determine the ambiguous effect on wages and employment, and the extent to which \glsxtrshort{LLM}'s create new tasks for a given occupation. First, \cite{noy} provide experimental evidence highlighting that ChatGPT benefits low skill workers by restructuring tasks and compressing the productivity distribution. \cite{hui2024short} use \glsxtrshort{DID} to estimate the effect of ChatGPT's release and other \glsxtrshort{LLM}'s on wages and employment for a large digital platform for freelancers. Their results suggest a fall in employment and wages for writers, designers and artists; especially those who are top rated. This last result may point to the possibility of shortening the wage gap between workers in the same exposed field. \cite{Dang} use the Current Population Survey (CPS) and a Synthetic \glsxtrshort{DID} and estimate the impact of ChatGPT's release on the united states's labour market. They find opposite results to \cite{hui2024short} in the form of a modest increase in wages for the exposed group and relatively no changes in unemployment for all population. \cite{humlum} rejects the idea that ChatGPT has a disruptive effect on the danish labor market and finds no effects on wages or employment. However, the authors do point out a major reorganization of tasks in exposed occupations. In summary, the literature related to \glsxtrshort{LLM} and its effects on the labor market does not follow a clear pattern.

Finally, the developing country evidence highlights certain dynamics developing countries do not face. Phenomenons such as informality complicate certain aspects of potential effects of \glsxtrshort{LLM}. Moreover, \cite{neto} nuances the job polarization narrative in developed economies for developing economies including Latin America. The authors find a moderated job polarization that is explained by limited technology adoption, structural change and changes in the global value chains. On the other hand, exposure to \glsxtrshort{AI} shows a clear pattern in Latin America.  Higher exposure is concentrated among women, urban, educated and formal workers \citep{azuara,worldbank,martinez2024risk}. This implies a major concern for labor inequality in the region in the following years. Specifically, \cite{martinez2024risk} adapt methodology by \cite{eloundou2023gptsgptsearlylook} and replicate their exposure index for the mexican context with \glsxtrfull{SINCO}\footnote{ \glsxtrshort{SINCO} enlists around 3,000 tasks for 468 occupations}. I use their measures in order to characterize occupations potentially exposed across several groups. The effort to estimate the ambiguous effect on employment and wages is uncharted territory in Mexico. My analysis is the first empirical approach to labor market outcomes in Mexico and one of the first for developing economies.




